% ============================================================
% 08_discusion.tex
% ============================================================

\section{Discusión}
\label{sec:discusion}

\subsection{Decisiones de diseño y sus fundamentos}

\subsubsection{Lógica de control en el firmware vs. en el servidor}

Una decisión de diseño central en NodeAlert IoT es la ubicación de la lógica de
control de iluminación íntegramente en el firmware del nodo sensor, sin depender de
instrucciones provenientes del servidor. Este enfoque, conocido como
\textit{edge computing}, ofrece ventajas sustanciales en entornos IoT donde la
conectividad de red puede ser intermitente: el sistema continúa operando
correctamente durante cortes de Wi-Fi o del broker MQTT, garantizando que la
iluminación del aula responda a las condiciones reales de ocupación y luminosidad
independientemente del estado de la red.

La contrapartida es que los parámetros de control (umbrales de luminosidad,
tiempo de retardo de apagado) son constantes de compilación, lo que requiere
reflashear el firmware para modificarlos. Esta limitación fue reconocida en el
diseño y se documentó explícitamente en la API (código HTTP 405 para
\texttt{POST /api/nodes/:id/config}).

\subsubsection{Separación entre sensor de ventana y sensor de aula}

El uso de dos sensores BH1750 en posiciones físicas distintas (uno orientado hacia
la fuente de luz natural, otro midiendo la iluminación interna del aula) resuelve
un problema inherente al control de iluminación artificial: el sensor interno
registra el efecto del propio foco encendido, creando un lazo de retroalimentación
positiva que podría mantener el foco apagado aun cuando la luz natural fuese
insuficiente. La lógica implementada en \texttt{regla\_deseado()} contempla esta
interacción: el sensor de ventana es el que determina si se enciende la luz,
mientras que el sensor de aula solo interviene cuando el sensor de ventana no está
disponible o cuando se requiere confirmar que el aula no tiene ya suficiente
iluminación por otro medio.

\subsubsection{Radar de microondas vs. sensor PIR}

La selección del HLK-LD2420 en lugar de un sensor PIR convencional se justifica
por dos razones técnicas. Primero, los sensores PIR son insensibles a personas
estáticas (sentadas en un escritorio trabajando), mientras que el radar FMCW
detecta la microvibración respiratoria. Segundo, los sensores PIR son fuertemente
afectados por corrientes de aire y cambios de temperatura ambiental. Sin embargo,
el radar FMCW implica mayor complejidad de integración: requiere protocolo UART
propio y un parser de texto no trivial, además de una lógica de votación para
reducir falsos positivos.

\subsubsection{Contenedores Docker vs. instalación nativa en Raspberry Pi}

El uso de Docker Compose para el despliegue del stack servidor aporta portabilidad,
reproducibilidad y aislamiento entre servicios. El comando de despliegue es una
única instrucción: \texttt{docker compose up -d --build}. La interdependencia entre
servicios se gestiona mediante el mecanismo de \textit{healthcheck} de Docker:
Node-RED no arranca hasta que el contenedor Mosquitto responda correctamente en el
puerto 1883.

\subsubsection{Node-RED como backend}

La elección de Node-RED como motor de backend permite implementar toda la lógica del
servidor —ingesta MQTT, persistencia SQLite, API REST, detector de nodos offline y
limpieza periódica— en un entorno visual sin necesidad de un servidor web
independiente. El soporte nativo de HTTP en Node-RED (nodos \textit{http in} y
\textit{http response}) evita la complejidad de integrar un framework web
adicional. El script \texttt{generate-flows.js} garantiza que los flujos sean
generados programáticamente y reproduzcan fielmente el diseño documentado, lo que
facilita el control de versiones y la modificación sistemática.

\subsection{Comparación con trabajos relacionados}

NodeAlert IoT comparte principios de diseño con el sistema descrito por
\textcite{Jia2019}, que emplea sensores de ocupación e iluminación para la gestión
energética en edificios educativos mediante protocolos IoT estándar. La principal
diferencia radica en el uso del radar FMCW en lugar de sensores PIR, y en la
implementación de toda la infraestructura sobre hardware de bajo costo (ESP32-C3
y Raspberry Pi) con software de código abierto.

Respecto a la arquitectura de middleware descrita por \textcite{Cardiel2012}, el
presente sistema adopta el modelo publicador/suscriptor de MQTT como capa de
comunicación entre el nodo sensor y el servidor de procesamiento, análogo al
concepto de middleware de adquisición de datos SCADA. La capa de lógica de negocio
(Node-RED) gestiona la persistencia, la exposición de datos y la comunicación hacia
la interfaz de usuario, en correspondencia con la separación de responsabilidades
propuesta por dichos autores.

\subsection{Limitaciones identificadas}

Las principales limitaciones del sistema actual son:

\begin{enumerate}
  \item \textbf{Umbrales de control fijos}: los parámetros de iluminación
        (60\,lx y 100\,lx) y el tiempo de retardo de apagado (300\,s) son
        constantes de compilación. No pueden modificarse en tiempo de ejecución
        sin reflashear el firmware.
  \item \textbf{Nodo único implementado}: el prototipo comprende un solo nodo
        sensor (\texttt{nodealert-luz01}). La arquitectura del sistema es
        multi-nodo por diseño (el tópico MQTT usa comodín \texttt{+}), pero la
        validación se realizó con un único dispositivo físico.
  \item \textbf{Sin encriptación TLS en MQTT}: la comunicación entre el ESP32-C3
        y el broker se realiza en texto plano sobre el puerto 1883. Para un entorno
        de producción, se requeriría TLS sobre el puerto 8883 con certificados.
  \item \textbf{Ausencia de historial gráfico de iluminación}: la vista de
        NodeDetail presenta gráficos de temperatura, gas y llama, pero no de
        iluminancia ni estado del relé. Estos datos están disponibles en la base
        de datos pero no se visualizan en gráficos temporales.
  \item \textbf{Sin calibración de sensores}: los valores reportados por los
        sensores BH1750 y SHT30 no han sido calibrados contra instrumentos de
        referencia; se utilizan los valores nominales de los datasheets de los
        fabricantes.
\end{enumerate}

